%% ma: L11-B14-0023
%% lop: 11
%% chuong: 4
%% bai: 14
%% dang: 11.14.1
%% loai: DS
%% muc: [TH, NB, TH, TH]
%% dap_an: "SĐSĐ"
%% nguon: "tai-lieu-hoc-tap-mon-toan-lop-11-theo-chuong-trinh-giao-duc-pho-thong-2018.pdf (Lớp Toán Thầy Duy, Phạm Lê Duy, THPT Chu Văn An), Chương 4 Bài 5, Phần 2, Câu 3 (tr.532)"
%% kien_thuc: "Lớp 11 Bài 14 (phép chiếu song song)"
%% trang_thai: chinh-thuc
%% ghi_chu: "Đã duyệt 03/11/2026."
\begin{ex}
Cho hình hộp chữ nhật $ABCD.A^{\prime}B^{\prime}C^{\prime}D^{\prime}$. Xét tính đúng, sai các mệnh đề sau:
\choiceTF
{Hình chiếu của $C$ lên $(A^{\prime}B^{\prime}C^{\prime})$ theo phương chiếu $DA^{\prime}$ trong phép chiếu song song là điểm $C^{\prime}$.}
{\True Gọi $O$ và $O^{\prime}$ lần lượt là tâm của $ABCD$ và $A^{\prime}B^{\prime}C^{\prime}D^{\prime}$. Hình chiếu của điểm $A$ lên mặt phẳng $(A^{\prime}B^{\prime}C^{\prime})$ theo phương chiếu $OO^{\prime}$ trong phép chiếu song song là điểm $A^{\prime}$.}
{Gọi $I$ và $I^{\prime}$ lần lượt là trung điểm của $AB$ và $A^{\prime}B^{\prime}$. Hình chiếu của $I$ lên mặt phẳng $(A^{\prime}B^{\prime}C^{\prime})$ theo phương chiếu $AI^{\prime}$ trong phép chiếu song song là điểm $A^{\prime}$.}
{\True Gọi $G$ là trọng tâm của $\triangle ABC$. Gọi $G^{\prime}$ là hình chiếu của $G$ lên $(A^{\prime}B^{\prime}C^{\prime}D^{\prime})$ theo phương chiếu $II^{\prime}$ trong phép chiếu song song. Khi đó $G^{\prime}$ là trọng tâm $\triangle A^{\prime}B^{\prime}C^{\prime}$.}
\loigiai{\sloppy a) $DA^{\prime}\parallel CB^{\prime}$ nên đường thẳng qua $C$ song song $DA^{\prime}$ là $CB^{\prime}$, cắt $(A^{\prime}B^{\prime}C^{\prime})$ tại $B^{\prime}$. Hình chiếu là $B^{\prime}$, không phải $C^{\prime}$. Sai.
b) $OO^{\prime}\parallel AA^{\prime}$ nên $A\mapsto A^{\prime}$. Đúng.
c) $\overrightarrow{AI^{\prime}}=\overrightarrow{AA^{\prime}}+\dfrac12\overrightarrow{AB}$. Điểm $I+\overrightarrow{AI^{\prime}}=B^{\prime}$ thuộc $(A^{\prime}B^{\prime}C^{\prime})$ nên hình chiếu của $I$ là $B^{\prime}$, không phải $A^{\prime}$. Sai.
d) $II^{\prime}\parallel AA^{\prime}$ nên $A,B,C\mapsto A^{\prime},B^{\prime},C^{\prime}$, và trọng tâm $G$ của $\triangle ABC$ có hình chiếu là trọng tâm của $\triangle A^{\prime}B^{\prime}C^{\prime}$. Đúng.}
\end{ex}
